\documentclass[11pt]{article}
\usepackage[a4paper,margin=26mm]{geometry}
\usepackage{amsmath,amssymb,mathtools}
\usepackage{unicode-math}
\setmainfont{Latin Modern Roman}
\setmathfont{Latin Modern Math}
\usepackage{xurl,hyperref}
\hypersetup{hidelinks,pdftitle={Problem statement}}
\newcommand{\sref}[2]{\hyperlink{source:#1:#2}{[S#2]}}
\newcommand{\cataloguescope}[1]{#1}
\begin{document}
\section{Numerical invariants — Question 13.7}\label{statement}

\subsection{Definitions and mathematical statement}

\paragraph{Shared notation.}$\mathbb N=\{1,2,\ldots\}$, and $\mathbb Z,\mathbb Q,\mathbb R,\mathbb C$ denote the integers, rationals, reals and complex numbers. Any other conventions are specified locally.


A surface lamination $\Lambda$ is a closed subset of $M$ locally homeomorphic to a product of a disk with a closed transverse set. It is essential when it has no sphere leaf or torus leaf bounding a solid torus, and its complementary regions are irreducible, have incompressible boundary, and have no compressing monogon. Such a monogon is a compressing disk with one boundary arc on a boundary leaf and the other along a transverse interval fiber. It is genuine if some complementary region is not an interval bundle. A sublamination is a closed union of leaves. For a cooriented codimension-one foliation choose a nowhere-zero $1$-form $\alpha$ with $T\mathcal F=\ker\alpha$. Integrability gives $d\alpha=\alpha\wedge\omega$. The Godbillon--Vey class is $\operatorname{gv}(\mathcal F)=[\omega\wedge d\omega]\in H^3(M;\mathbb R)$, independent of these choices, and its value on an oriented closed $M$ is the integral of this $3$-form.
\paragraph{Statement recorded in the supplied source.}
Construct a suitable Godbillon--Vey-type invariant for laminations, specifying its domain, definition and invariance properties. The source includes a weaker possible direction: a well-defined positive, negative or zero twisting invariant according to whether suitable geometric realizations can make transverse helical twisting have one sign. The requested construction should account for complementary regions rather than assume a lamination is already a foliation. \sref{statement}{2}
\subsection{Short English statement}
Develop a Godbillon--Vey-type invariant for laminations, possibly beginning with a well-defined sign of transverse twisting.
\subsection{Sources}
\begin{description}
\item[\hypertarget{source:statement:1}{S1}] ULAM.AI and the UnsolvedMath Contributors, \emph{UnsolvedMath: A Curated Collection of Open Mathematics Problems}, catalogue entry. CC BY 4.0. \url{https://huggingface.co/datasets/ulamai/UnsolvedMath}.
\item[\hypertarget{source:statement:2}{S2}] Danny Calegari, \emph{Problems in foliations and laminations of 3-manifolds}, version 0.78 (2002), Question 13.7; Definitions 1.1--1.10 and local definitions and remarks. \url{https://arxiv.org/abs/math/0209081v1}.
\end{description}
\label{end:statement}
\end{document}
